% !TEX root = ../notes.tex

\documentclass[../notes.tex]{subfiles}

\begin{document}

\section{April 6}
Before officially starting, let's say something about animation.
\begin{remark}[animation]
	The category of connective $\ZZ$-modules is equivalent to the category of simplicial abelian groups, considered up to $\pi_*$-isomorphisms. (This is a version of the Dold--Kan correspondence; the point is that $\pi_*$-isomorphisms go to equivalences of the $\ZZ$-modules.) By so-called ``animation,'' this turns out to be equivalent to the category of product-preserving functors from finitely generated free abelian groups to $\mathrm{Spaces}$; roughly speaking, we are freely adding ``sifted'' colimits. Similarly, for any discrete ring $R$, the category of connective $R$-modules is equivalent to the product-preserving functors from free finitely generated $R$-modules to spaces.
\end{remark}
The reason we are bringing up animation is that a formal group over a connective $\mathbb E_\infty$-ring $R$ is a product-preserving functor from finitely generated free abelian groups to formal hyperplanes over $R$. Note that we need formal hyperplanes in more than one dimension to even make sense of produces in the target! Note that formal hyperplanes can be embedded into the category of functors from connective $\mathbb E_\infty$-algebras to spaces, so we are more or less describing one of these animation functors which have the extra structure of a formal hyperplane.

\subsection{Where to Find Formal Groups}
Here is our definition.
\begin{definition}[weakly even periodic]
	An $\mathbb E_\infty$-ring is \textit{weakly even periodic} if and only if $\pi_*E$ is concentrated in even degrees, and the maps
	\[\pi_2E\otimes_{\pi_0E}\pi_iE\to\pi_{i+2}E\]
	are isomorphisms for each $i$.
\end{definition}
\begin{remark}
	There does not need to be a unit in $\pi_2E$, which is typically how an $\mathbb E_\infty$-ring becomes weakly even periodic. Nonetheless, $\pi_2E$ should be rank $1$ and projective by taking $i=-2$.
\end{remark}
\begin{lemma} \label{lem:get-smooth-coalg}
	Suppose that $E$ is weakly even periodic. Then $E\otimes\Sigma_+^\infty\CP^\infty$ is a smooth coalgebra.
\end{lemma}
\begin{proof}
	Indeed,
	\[E\otimes\Sigma_+^\infty\CP^n=E\oplus\Sigma^2\oplus\cdots\oplus\Sigma^{2n}E\]
	by some cell decomposition: this is true for $N=0$ with no content, and the general case follows by using the following pushout square.
	% https://q.uiver.app/#q=WzAsNCxbMCwxLCJcXFNpZ21hXlxcaW5mdHlfK1xcbWF0aHJte3B0fSJdLFsxLDEsIlxcU2lnbWFfK15cXGluZnR5IFxcQ1BebiJdLFswLDAsIlxcU2lnbWFeXFxpbmZ0eV8rU157Mm4rMX0iXSxbMSwwLCJcXFNpZ21hXlxcaW5mdHlfK1xcQ1Bee24tMX0iXSxbMiwzXSxbMywxXSxbMCwxXSxbMiwwXV0=&macro_url=https%3A%2F%2Fraw.githubusercontent.com%2FdFoiler%2Fnotes%2Fmaster%2Fnir.tex
	\[\begin{tikzcd}[cramped]
		{\Sigma^\infty_+S^{2n+1}} & {\Sigma^\infty_+\CP^{n-1}} \\
		{\Sigma^\infty_+\mathrm{pt}} & {\Sigma_+^\infty \CP^n}
		\arrow[from=1-1, to=1-2]
		\arrow[from=1-1, to=2-1]
		\arrow[from=1-2, to=2-2]
		\arrow[from=2-1, to=2-2]
	\end{tikzcd}\]
	Then we take $E\otimes-$ and note that the top horizontal arrow is $0$ (by induction and the fact that $E$ is even), so we have a fiber sequence, and the claim follows. It follows that $\pi_*\left(E\otimes\Sigma_+^\infty\CP^\infty\right)=\bigoplus_{n\ge0}\pi_{*+2n}E$. It now follows that $E\otimes\Sigma_+^\infty\CP^\infty$ is a smooth coalgebra by staring at the definition.
\end{proof}
Thus, we have access to many formal hyperplanes.
\begin{example}
	By \Cref{lem:get-smooth-coalg}, we see that $\tau_{\ge0}\left(E\otimes\Sigma_+^\infty\CP^\infty\right)$ is also a smooth coalgebra over $\tau_{\ge0}E$, so
	\[\op{Spf}\op{Hom}_{\tau_{\ge0}E}\left(\tau_{\ge0}\left(E\otimes\Sigma_+^\infty\CP^\infty\right),\tau_{\ge0}E\right)\]
	is a formal hyperplane.
\end{example}
\begin{example}
	The Quillen formal group over $\tau_{\ge0}E$ is the functor
	\[M\mapsto\op{Spf}\op{Hom}_{\tau_{\ge0}E}\left(\tau_{\ge0}(E\otimes\Sigma_+^\infty K(M^\lor,2)),\tau_{\ge0}E\right)\]
	from finitely generated free abelian groups to formal hyperplanes. The moral is that we always get a formal group from our rings, in the same way that $E^0\CP^\infty$ is a classical formal group over $\pi_0E$.
\end{example}
The above example motivates why Lurie wants formal groups to lift to $\ZZ$-modules: this is the structure which we already have from $E^0\CP^\infty$.

\subsection{\texorpdfstring{$p$}{ p}-divisible Groups}
By way of example, note that the classical group $\mathbb G_m$ admits a formal group
\[\widehat{\mathbb G}_m=\colim\Spec\frac{\ZZ[t,1/t]}{\left(t-1\right)^\bullet}\]
and also the $p$-divisible group
\[\mu_{p^\infty}=\colim\Spec\frac{\ZZ[t]}{\left(t^\bullet-1\right)}.\]
In general, these things are rather different, but they agree when $p$ is nilpotent. Roughly speaking, we can start our lives studying formal groups, but $p$-divisible groups have a tendency to be more concrete (e.g., from the perspective of number theory). Here is our definition in homotopy theory.
\begin{defihelper}[$p$-divisible group] \nirindex{p-divisible group@$p$-divisible group}
	Fix a connective $\mathbb E_\infty$-ring $E$, and choose a prime $p$. Then a \textit{$p$-divis\-ible group} over $R$ is a functor $\mathbb G$ from connective $\mathbb E_\infty$-algebras to connective $\ZZ$-modules satisfying the following.
	\begin{itemize}
		\item For all $A$, we have $\mathbb G(A)[1/p]=0$.
		\item The functor
		\[A\mapsto\Omega^\infty\op{Hom}_\ZZ(\ZZ/p\ZZ,\mathbb G(A))\]
		is representable by a finite flat $R$-algebra.
		\item The map $p\colon\mathbb G\to\mathbb G$ is locally surjective in the finite flat topology. Explicitly, for any $A$ and $x\in\pi_0\mathbb G(A)$, there is a finite flat map $A\to B$ of connective $\mathbb E_\infty$-algebras (which requires that $\pi_0A\to\pi_0B$ is flat, and $\pi_iA\otimes\pi_jB\cong\pi_{i+j}B$) such that $\pi_0A\to\pi_0B$ is faithfully flat, and the image of $x$ in $\pi_0\mathbb G(B)$ is divisible by $p$.
	\end{itemize}
\end{defihelper}
Roughly speaking, by an infinitesimal approximation, one can recover a formal group from its $p$-divisible group.
\begin{theorem}
	Fix a $p$-complete connective $\mathbb E_\infty$-ring $R$. For any $p$-divisible group $\mathbb G$, there is a unique formal group $\mathbb G^\circ$ over $R$ such that the following holds: if $A$ is a connective $\mathbb E_\infty$-algebra such that $p$ is nilpotent in $\pi_0A$ and $\pi_iA=0$ for $i$ large, then there is a fiber sequence
	\[\mathbb G^0A\to\mathbb G(A)\to\mathbb G(\pi_0A^{\mathrm{red}}).\]
\end{theorem}
One may hope that the reduction doesn't do anything.
\begin{definition}[connected]
	Fix a $p$-complete $\mathbb E_\infty$-ring $R$. Then a $p$-divisible group $G$ over $R$ is \textit{connected} if and only if the following holds: if $A$ is a connective $\mathbb E_\infty$-algebra such that $p$ is nilpotent in $\pi_0A$ and $\pi_iA=0$ for $i$ large, then $\mathbb G(A)=\mathbb G^\circ(A)$.
\end{definition}
\begin{example} \label{ex:mu-p-connected}
	We claim that $\mu_{p^\infty}$ is connected. Indeed, for this, we should check that $\mu_{p^\infty}(A)=1$ if $p$ is nilpotent in $A$ and $A$ is reduced (which in particular means that $pA=0$).
\end{example}
\begin{remark}
	Thus, over a $p$-complete ring, the category of connected $p$-divisible groups embeds into the category of formal groups.
\end{remark}
We have now worked hard enough to earn some examples.
\begin{example} \label{ex:etale-p-div}
	The constant group scheme $\QQ_p/\ZZ_p$ is a non-connected $p$-divisible group over $\mathbb S_p$. Explicitly, this functor takes $p$-complete connective $A$ to the locally constant function $\Spec\pi_0A\to\QQ_p/\ZZ_p$. In particular, this has interesting values on reduced rings!
\end{example}
\begin{remark}
	Any $p$-divisible group is an extension of an \'etale $p$-divisible group by a connected one. Connected $p$-divisible groups look like \Cref{ex:mu-p-connected}, and \'etale $p$-divisible groups look like \Cref{ex:etale-p-div}.
\end{remark}
Now, one may hope that we can take a formal group and try to take torsion or similar to recover a $p$-divisible group, but we cannot always do so.
\begin{exe}
	The formal group $\widehat{\mathbb G}_a$ is not $p$-divisible over $\FF_p$.
\end{exe}
\begin{proof}
	Let's state the problem we are trying to solve: given a formal group $\widehat{\mathbb G}$, we are hunting for a $p$-divisible group $\mathbb G$ for which $\widehat{\mathbb G}=\mathbb G^\circ$, meaning that the functors
	\[A\mapsto\Omega^\infty\op{Hom}_\ZZ(\ZZ/p\ZZ,\mathbb G(A))\]
	and
	\[A\mapsto\op{fiber}\left(p\colon\widehat{\mathbb G}(A)\to\widehat{\mathbb G}(A)\right)\]
	and be co-represented by a finite flat $R$-algebra. In our example, we are hunting for a formal group law $F$ over $\FF_p$, which is some addition law for $\FF_p[[x]]$. Taking $\pi_0$ and so on, one translates that we need $\FF_p[[x]]/(x+_F\cdots+_Fx)$ needs to be finite flat, but it is not for $\widehat{\mathbb G}_a$: this quotient is $\FF_p[[x]]/(px)=\FF_p[[x]]$!
\end{proof}
\begin{remark}
	For $\widehat{\mathbb G}_m$, we can see that $\FF_p[[x]]/([p]_Fx)$ is in fact finite: this quotient is finite-dimensional over $\FF_p$!
\end{remark}
We can now state our target theorem.
\begin{definition}[deformation]
	Fix a $p$-divisible group $\mathbb G$ over a classical commutative ring $R_0$. Then a \textit{deformation} of $\mathbb G_0$ to $R$ is the data of a map $R\to R_0$ along with a $p$-divisible group $\mathbb G$ over $R$ equipped with an isomorphism $\mathbb G_{R_0}\cong\mathbb G_0$. We let $\op{Def}_{G_0}(-)$ denote the corresponding functor from $R$-algebras to spaces, and we let $\op{Def}_{\mathbb G_0}(-)^\simeq$ be the groupoid.
\end{definition}
\begin{theorem}
	Fix a $p$-divisible group $\mathbb G_0$ over a classical ring $R_0$. Suppose that $p$ is nilpotent in $R_0$, that $L_{R_0}$ is almost perfect as an $R_0$-module, and $\mathbb G_0$ is non-stationary. Then the functor $\op{Def}_{\mathbb G_0}(-)^\simeq$ is co-representable by an $\mathbb E_\infty$-ring.
\end{theorem}
\begin{remark}
	If $R_0$ is perfect over $\FF_p$, then all these conditions are automatic.
\end{remark}

\end{document}